\documentclass[10pt,a4paper]{amsart}
\usepackage[margin=19mm]{geometry}
\usepackage{amsmath,amssymb,mathtools}
\usepackage[scr=rsfs,cal=euler]{mathalfa}
\usepackage{xfrac}
\usepackage[unicode,hidelinks,bookmarks=false]{hyperref}
\newcommand{\ie}{i.e.\ }
\newcommand{\cf}{cf.\ }
\newcommand{\resp}{respectively\ }
\newcommand{\Subsection}{\subsection}
\providecommand{\nobreakdash}{\nobreak\hbox{-}}
\setlength{\emergencystretch}{2em}
\multlinegap0pt
\usepackage{amssymb}
\newtheorem{theorem}{Theorem}
\newcommand{\R}{\mathbb{R}} 	%real numbers
\title{Helly-type problems from a topological perspective}
\author{Pavel Pat\'ak, Zuzana Pat\'akov\'a}
\date{Source: arXiv:2602.07552v1 (CC BY 4.0)}
\begin{document}
\maketitle

\noindent\emph{Source and licence.} Helly-type problems from a topological perspective. Version arXiv:2602.07552v1 (CC BY 4.0), distributed by arXiv under the Creative Commons Attribution 4.0 International licence, http://creativecommons.org/licenses/by/4.0/, as recorded in the arXiv licence metadata for this exact version and shown on the abstract page https://arxiv.org/abs/2602.07552v1. Full licence notice: \texttt{licenses/arxiv-2602.07552-CC-BY-4.0.txt}.\par\smallskip

\noindent\textit{Editorial scope.} Counted unit: the article's theorem thm:helly (source0.tex lines 84--85). The article's complete original proof of it is delivered with it (source0.tex lines 331--339, one \texttt{proof} environment of 9 lines, which the article places away from the statement). No mathematics is written, repaired or re-derived: the statement and the proof are the article's own lines, in the article's own notation, and no retained line is substituted or re-flowed.  No further source-local context is needed: every cross-reference of every retained line resolves inside the two delivered spans.  Nothing was omitted from any retained span, so the package carries no omission record.  The retained lines cite 1 works, whose entries are transcribed verbatim from the article's own bibliography source in the e-print and delivered with short editorial labels recorded in the item's reference mapping.  No retained line contains a figure environment, an image-inclusion command or drawing code, so no image is delivered and the package declares no asset dependency.  Editorial formatting only, in addition: an AMS \texttt{amsart} preamble that restates the article's own declarations and macros used by the retained lines, namely the environments \texttt{theorem} on separate counters, together with the standard packages those lines need.  The retained environments print in this document as Theorem 1 in order of appearance, whereas the article prints the same items with the numbering of its own full document.  The complete native source of the article as distributed in the arXiv e-print for this exact version is bundled unchanged as \texttt{sources/194187/source0.tex}.
\begin{theorem}[Helly's theorem~\cite{h-umkkm-23}]\label{thm:helly} Given a finite family $\mathcal F$ of convex sets in $\R^d$, either $\bigcap \mathcal F\neq \emptyset$, or one can find a subfamily $\mathcal G\subseteq \mathcal F$ that has at most $d+1$ elements and satisfies $\bigcap \mathcal G=\emptyset$.
\end{theorem}

\begin{proof}[Proof of Helly's theorem (Thm.~\ref{thm:helly}) from Radon's lemma]
For contradiction assume that the Helly number of convex sets in $\R^d$ is larger than $d+1$. Therefore, there are $n>d+1$ convex sets $G_1,G_2,\ldots, G_n$ for which $G_1\cap\cdots\cap G_n=\emptyset$ and for all $i=1,\ldots, n$, $\bigcap_{j\neq i}G_j\neq\emptyset$. Therefore, for each $i=1,\ldots, n$, we may choose a point $p_i$ that lies in $\bigcap_{j\neq i}G_j$.
Let $f\colon \Delta_{n-1}\to \R^d$ be the linear map that maps the vertex $e_i$ of $\Delta_{n-1}$ to the point $p_i$, $i=1,\ldots, n$.
Then for each non-empty face $\sigma$ of $\Delta_{n-1}$
$f(\sigma)\subseteq \bigcap_{e_j\notin \sigma}G_j$.
Since $n\geq d+2$, Radon's lemma yields two disjoint faces $\sigma,\tau$ of $\Delta_{n-1}$ for which
$f(\sigma)\cap f(\tau)\neq\emptyset$.
But then, the intersection is contained in $G_1\cap \cdots\cap G_n$, which is thus non-empty, a contradiction.
\end{proof}

\begin{thebibliography}{99}
\bibitem[h-umkk]{h-umkkm-23} E.~Helly. \newblock {\"U}ber {Mengen} konvexer {K{\"o}rper} mit gemeinschaftlichen {Punkten}. \newblock {\em Jahresbericht Deutsch. Math. Verein.}, 32:175--176, 1923.
\end{thebibliography}
\end{document}
